\begin{textsample}{NEG Dim 8 – human – Score -2.00 – rj/rj\_ef\_ciencias\_da\_natureza\_2\_p7.txt}  \label{ex:f8_neg_005}
Nos anos finais, a ampliação da relação dos estudantes com o ambiente possibilita expandir a exploração dos fenômenos relacionados aos materiais e energia, no âmbito do sistema produtivo e seu impacto na qualidade ambiental.
Assim, o aprofundamento da temática dessa unidade, que envolve inclusive a construção de modelos explicativos, deve possibilitar aos estudantes fundamentar -se no conhecimento científico para, por \textbf{exemplo} : \textbf{avaliar} vantagens e desvantagens da produção de produtos sintéticos a partir de recursos naturais, da produção e do uso de determinados combustíveis, bem como da produção, da transformação e da propagação de diferentes tipos de energia e do funcionamento de artefatos e equipamentos que propiciam novas formas de interação com o ambiente, estimulando tanto a reflexão para hábitos mais sustentáveis no uso dos recursos naturais e científico-tecnológicos quanto à produção de novas tecnologias e o desenvolvimento de ações coletivas de aproveitamento responsável dos recursos.

% matched lemmas: avaliar, exemplo
\end{textsample}
